\section{Complex Polynomials}
\label{ch:ancora_polinomi}

\subsection{The Fundamental Theorem of Algebra}

The fundamental theorem of algebra states that every non-constant polynomial has
at least one zero in the complex plane.
\mynote{see the historical notes at the end of the chapter}
To prove the theorem, we must extend Weierstrass's theorem to functions of a
complex variable.
In the real Weierstrass theorem, the function is defined on a closed and bounded
interval.
In the complex plane, the concept of an \emph{interval} does not exist since
there is no ordering, but we will see that Weierstrass's theorem remains valid
for continuous functions defined on closed and bounded sets according to the
following definitions.

\begin{definition}[Sequential Closure]
A set $A\subset \CC$ is said to be \emph{sequentially closed}
\mymargin{sequentially closed}%
\index{sequentially!closed}%
if for any sequence of points $a_n\in A$, if $a_n \to a$ for some $a\in \CC$,
then $a\in A$.
\end{definition}

\begin{definition}[Boundedness]
A set $A\subset \CC$ is said to be \emph{bounded}%
\mymargin{bounded}%
\index{bounded}
if
\[
  \sup \ENCLOSE{ \abs{z}\colon z \in A} < +\infty.
\]
\end{definition}

Note that a sequence $a_n\in \CC$ is bounded (see
chapter~\ref{sec:successione_limitata}) if and only if its image
$\ENCLOSE{a_n\colon n\in \NN}\subset \CC$ is a bounded set.

\begin{theorem}[Complex Weierstrass]
Let $A\subset \CC$ be a non-empty, sequentially closed, and bounded set, and let
$f\colon A \to \RR$ be a continuous function.
Then $f$ attains a maximum and a minimum on $A$.
\end{theorem}
%
\begin{proof}
We prove the existence of the minimum; the proof for the maximum is perfectly
analogous.
Let $m=\inf f(A)$.
Since $A$ is non-empty, by the lemma on the existence of minimizing sequences,
there exists a sequence $z_n \in A$ such that $f(z_n) \to m$.
Since $A$ is bounded, we can apply the Bolzano-Weierstrass theorem to find $z\in
\CC$ and a subsequence $z_{n_k} \to z$. Since $A$ is sequentially closed, we can
assert that $z\in A$. Since $f$ is continuous, we conclude that
\[
f(z) = \lim_{k\to+\infty} f(z_{n_k}) = m
\]
and thus $z$ is a point of minimum for $f$.
\end{proof}

\begin{theorem}[Existence of Minimum for Coercive Functions]
Let $f\colon \CC \to \RR$ be a continuous function such that for every sequence
$z_n \to \infty$ (i.e., $\abs{z_n}\to +\infty$), we have $f(z_n) \to +\infty$.
Then $f$ attains a minimum.
\end{theorem}
%
\begin{proof}
Consider the set
\[
  A = \ENCLOSE{z \in \CC \colon f(z) \le f(0)}.
\]
Clearly, $0\in A$, so $A$ is non-empty.
The set $A$ is also sequentially closed because if $z_k\in A$, then $f(0) -
f(z_k)\ge 0$, and by continuity, $f(0)-f(z_k)\to f(0)-f(z)$, and by the sign
permanence theorem, we obtain $f(0)-f(z) \ge 0$, i.e., $z \in A$.
Now, we demonstrate that $A$ is also bounded.
If it were not, there would exist, by contradiction, a sequence $z_n \in A$ such
that $\abs{z_n}\to +\infty$, i.e., $z_n \to \infty$.
But then, by hypothesis on $f$, we would have $f(z_n)\to +\infty$, which
contradicts the condition $f(z_n) \le f(0)$.
Since $A$ is non-empty, sequentially closed, and bounded, and since $f\colon A
\to \RR$ is continuous, we can apply the complex Weierstrass theorem to deduce
that $f$ attains a minimum on $A$ at some point $w \in A$.
But since $0\in A$, we certainly have $f(w)\le f(0)$, and for every $z\in \CC
\setminus A$, we have $f(z) > f(0)$ by the definition of $A$.
Thus, $w$ is the minimum of $f$ on the entire $\CC$, not just on $A$.
\end{proof}

\begin{theorem}[Fundamental Theorem of Algebra]
\label{th:fondamentale_algebra}
\mymargin{fundamental theorem of algebra}%
\index{fundamental theorem of algebra}%
Let $f(z)$ be a polynomial of degree $N\ge 1$ with complex coefficients:
\[
  f(z) = \sum_{j=0}^N a_j \cdot z^j
\]
with $a_j\in \CC$ for $j=0,\dots,N$ and $a_N \neq 0$.
Then there exists $w\in \CC$ such that $f(w) = 0$.
\end{theorem}
%
\begin{proof}
    First, observe that $\abs{f(z)}$ is coercive, meaning that if $z_n \to
  \infty$, then $\abs{f(z_n)}\to +\infty$.
  Indeed, we have
  \begin{align*}
    \abs{f(z_n)}
    &= \abs{\sum_{j=0}^N a_j z_n^j}
    = \abs{a_N z_n^N  + \sum_{j=0}^{N-1} a_j z_n^j}\\
    &= \abs{z_n}^N \cdot \abs{a_N + \sum_{j=0}^{N-1} \frac{a_j}{z_n^{N-j}}}
    \to +\infty
  \end{align*}
  if $z_n \to \infty$.
  
    We know that all polynomials are continuous functions since they are sums of
  products of continuous functions, and the modulus is also a continuous
  function, so $\abs{f(z)}$ is certainly a continuous function.
  
    Thus, we can apply the theorem on the existence of a minimum for coercive
  functions: there exists $w\in \CC$ such that $\abs{f(w)}$ is minimal.
  
  To conclude the theorem, it suffices to show that $f(w)=0$.
    The idea we want to develop is that complex polynomials, if they assume a
  value $f(w)$ at a point $w\in \CC$, then they also assume all values close to
  it because \emph{locally}, the polynomial resembles a power $z^n$, and the
  equation $z^n=c$ always has a solution, as we have already seen.
    Thus, near $w$, there will be points where $f$ assumes values whose modulus
  is less than $f(w)$, unless $f(w)=0$, in which case it is obviously impossible
  to have numbers with a modulus less than $0$.
  
  We can write the polynomial $f(z)$ in the form
  \[
    f(z) = \sum_{j=0}^N b_j (z-w)^j
  \]
    since the translation of a polynomial is still a polynomial of the same
  degree.
  Proving that $f(w)=0$ is therefore equivalent to proving that $b_0=0$.
  Suppose, for contradiction, that $b_0\neq 0$. 
    The behavior of the polynomial near the point $w$ is dominated by the lowest
  degree terms: some of these may be zero, but we can consider the first index
  $k>0$ for which $b_k\neq 0$. Then we can write:
  \[
      f(z) = b_0 + b_k(z-w)^k + \sum_{j=k+1}^N b_j (z-w)^j.
  \]
Now write $z-w$, $b_0$, and $b_k$ in exponential form:
$z - w = \rho e^{i\theta}$, $b_0 = r e^{i \alpha_0}$,
$b_k = r_k e^{i\alpha_k}$ 
and apply the triangle inequality
  \begin{align*}
    \abs{f(z)} 
    &= \abs{r e^{i\alpha} + r_k \rho^k e^{i(\theta k+\alpha_k)} 
    + \sum_{j=k+1}^N b_j \rho^j e^{i\theta j}} \\
    &\le \abs{r e^{i\alpha} + r_k \rho^k e^{i(\theta k + \alpha_k )}} 
    + \sum_{j=k+1}^N \abs{b_j} \rho^j.
  \end{align*}
To reach a contradiction, we want to show that there exists $z$ (i.e., there
exist $\rho$ and $\theta$) for which the value of the function is less in
modulus than $r = \abs{b_0} = \abs{f(w)}$.
First, if $\rho$ is sufficiently small, we can make the powers $\rho^j$ with
$j>k$ arbitrarily small: it suffices to choose $\rho\le 1$ and $\rho \le
r_k/(2\sum \abs{b_j})$ so that we have
\[
  \sum_{j=k+1}^N \abs{b_j} \rho^j
  \le \rho^{k+1} \sum_{j=k+1}^N \abs{b_j} \le r_k\frac{\rho^{k}}{2}.
\]
Regarding the term $r e^{i\alpha} + r_k \rho^k e^{i(\theta k + \alpha_k)}$, it
will suffice to choose $\theta = (\pi + \alpha - \alpha_k) /k$ so that the two
addends have opposite phases and the sum is destructive:
\[
  \abs{r e^{i\alpha} + r_k \rho^k e^{i(\theta k + \alpha_k)}}
  = \abs{e^{i\alpha}(r + r_k \rho^k e^{i\pi})}
  = r - r_k\rho^k.
\]
Ultimately, we have
\[
\abs{f(z)} = \abs{f(w+\rho e^{i\theta})} 
\le r - r_k\rho^k + r_k\frac{\rho^k}{2} 
< r = \abs{f(w)}
\]
which is a contradiction since $\abs{f(w)}$ was supposed to be the minimum
value.
\end{proof}
  
\subsection{Factorization of Polynomials}

\begin{theorem}[Factorization of Complex Polynomials]
\index{decomposition!of complex polynomials}%
\index{factorization!of complex polynomials}%
\index{polynomial!complex factorization}%
Let $p(z)$ be a non-zero polynomial. Then, given $n=\deg p$, there exist complex
numbers $z_1, z_2, \dots, z_n$ and a non-zero complex number $c$ such that
\begin{equation}\label{eq:34985}
  p(z) = c \prod_{k=1}^n (z-z_k).
\end{equation}
The $z_k$ are unique up to order, and $c$ is uniquely determined.
\end{theorem}
%
\begin{proof}
We prove the theorem by induction on $n=\deg p$. If $n=0$, the polynomial $p$ is
constant: $p(z) = c$. Recalling that a product of $n=0$ factors is equal to $1$,
we obtain the desired result.

Now let $p(z)$ be any polynomial of degree $n>0$. By the fundamental theorem of
algebra, we know there exists a complex number $z_n$ such that $p(z_n)=0$. By
Ruffini's theorem, we have
\[
  p(z) = (z-z_n) q(z)
\]
with $q$ being some polynomial of degree $n-1$. By the inductive hypothesis, we
can then assume there exist complex numbers $z_1, \dots, z_{n-1}$ and $c$ such
that
\[
   q(z) = c \prod_{k=1}^{n-1} (z-z_k)
\]
and the thesis follows.
\end{proof}

In the factorization~\eqref{eq:34985}, the \emph{roots} $z_k$ may repeat. If we
group together the factors corresponding to the same root, we obtain a
decomposition of the form
\begin{equation}\label{eq:358925}
P(z) = c \prod_{k=1}^m (z-z_k)^{p_k}
\end{equation}
with $z_1, \dots, z_m$ being distinct complex numbers (the roots of the
polynomial $P$) and $p_k$ being positive integers.
The exponent $p_k$ is called the \emph{multiplicity}%
\mymargin{multiplicity}%
\index{multiplicity} of the root $z_k$, and it holds that
\[
  p_1 + \dots + p_m = n.
\]
This equality is expressed by saying that a polynomial $P\in \CC[z]$ of degree
$n$ always has $n$ roots counted with their multiplicity.
The distinct roots are instead $m\le n$.

If $P\in \RR[x]$ is a polynomial with real coefficients, it is not guaranteed to
have roots: for example, the polynomial $x^2+1$ has no real roots since
$x^2+1>0$ as happens in any ordered field.
It may then be useful to consider $P$ as a polynomial in $\CC[x]$ with real
coefficients.

If $P\in \RR[x]$ is a polynomial with real coefficients
\[
  P = \sum_{k=0}^n a_k x^k, \qquad a_k\in \RR
\]
we can consider $P$ also as a polynomial in $\CC[x]$, since $a_k\in \RR\subset
\CC$.
The following theorem provides a criterion for distinguishing polynomials with
real coefficients within $\CC[x]$.

\begin{theorem}
\label{th:caratterizzazione_polinomi_reali}%
If $P\in \CC[x]$ is a polynomial with complex coefficients
\[
  P = \sum_{k=0}^n a_k x^k,\qquad a_k\in \CC
\]
and if there exists an infinite set $A\subset \RR$ such that for every $x\in A$,
we have $P(x)\in \RR$, then $a_k\in\RR$ for every $k=0,\dots,n$, and thus $P\in
\RR[x]$.

In particular, if the polynomial function associated with $P\in\CC[x]$ maps
$\RR$ to $\RR$, then $P$ is a polynomial with real coefficients.
\end{theorem}
%
\begin{proof}
Consider the polynomial:
\[
  Q = \sum_{k=0}^n (a_k-\bar a_k) x^k.
\]
Then for every $x\in A$, we have
\[
  Q(x) = \sum_{k=0}^n a_k x^k - \sum_{k=0}^n \bar a_k x^k
     = P(x) - \overline{P(x)} = P(x) - P(x) = 0
\]
since if $x\in \RR$, it follows that
$\overline{x^k} = {\bar x}^k = x^k$.
Since $A$ is infinite,
by the polynomial zero theorem
(theorem~\ref{th:annullamento_polinomi}),
we deduce that $Q=0$.
But then all the coefficients of $Q$ must
be zero, i.e., $\bar a_k = a_k$.
It follows that $a_k\in \RR$ for every $k=0,\dots,n$.
\end{proof}

Referring to the factorization of polynomials with complex coefficients, we can
also factor polynomials with real coefficients if we are content with having
quadratic factors instead of linear ones.

\begin{theorem}[Factorization of Real Polynomials]
  \label{th:fattorizzazione_polinomio_reale}%
If $P\in \RR[x]$ is a polynomial with real coefficients, we can write
\begin{equation}\label{eq:35549}
    P = a \cdot \prod_{k=1}^n (x-x_k)^{p_k} \cdot \prod_{j=1}^m (x^2 + \alpha_j
  x + \beta_j)^{q_j}
\end{equation}
where $a\in \RR$, $x_k\in \RR$, $p_k\in \NN\setminus\ENCLOSE{0}$,
$\alpha_j,\beta_j\in \RR$, $q_j\in \NN\setminus\ENCLOSE{0}$
with $\alpha_j^2 - 4 \beta_j < 0$
and
\[
  \sum_{k=1}^n p_k + 2 \sum_{j=1}^m q_j = \deg P.
\]

The factorization~\eqref{eq:35549} is unique up to the order of the factors.

The numbers $x_1,\dots,x_n$ are all the distinct real roots of the polynomial
$P$ with respective multiplicities $p_1,\dots,p_n$, while all the distinct
complex (non-real) roots will be $\mu_1,\dots, \mu_m$ and $\bar \mu_1, \dots,
\bar \mu_m$ with
\[
  x^2 + \alpha_j x + \beta_j = (x-\mu)\cdot(x-\bar \mu)
\]
from which
\[
  \alpha_j = -2\Re \mu_j, \qquad \beta_j=\abs{\mu_j}^2
\]
and $q_j$ will be the multiplicities of the roots $\mu_j$ and $\bar \mu_j$.
\end{theorem}
%
\begin{proof}
First, observe that if $P$ has real coefficients, we have
\[
  \overline{P(z)} = P(\bar z), \qquad \forall z\in \CC
\]
thus if $\mu$ is a root of $P$, then $\bar \mu$ is also a root of $P$.
Now, if $\mu$ is a non-real root of $P$, we know that in the complex field, $P$
is divisible by $x-\mu$ (Ruffini's theorem~\ref{th:Ruffini}):
\[
  P = (x-\mu) \cdot P_1.
\]
But $\bar \mu$ is also a root of $P$, and since $\bar \mu \neq \mu$, we must
have
\[
Q(\bar \mu) = \frac{P(\bar \mu)}{\bar \mu - \mu} = 0
\]
and thus, applying Ruffini's theorem again,
\[
 P = (x-\mu)\cdot (x-\bar \mu)\cdot P_2.
\]
Now observe that
\[
(x-\mu)\cdot(x-\bar \mu) = x^2 - (\mu + \bar \mu) x + \mu \bar \mu
 = x^2 - 2 (\Re \mu)\cdot x + \abs{\mu}^2
\]
is a polynomial with real coefficients, and since $P_2$ is obtained by dividing
$P$ by this polynomial, $P_2$ is also a polynomial with real coefficients.

Repeating the same procedure on the polynomial $P_2$, we can factor $P$ by
reducing the degree in steps of $2$ until all non-real complex roots are
exhausted by pairing them two by two.
Thus, if $\mu$ is a complex root of the real polynomial $P$, the multiplicity of
$\mu$ is equal to the multiplicity of $\bar \mu$.

Afterward, we can complete the factorization by dividing by the linear factors
$x-x_k$ corresponding to the real roots of the polynomial $P$.
We will then obtain the desired factorization.
\end{proof}